\section{Код в Wolfram Mathematica для решения уравнения осцилляции нейтрино в среде при разных значениях \(\xi\)}
\label{app: code-python-vis}

Для численного решения уравнения \eqref{eq: 15} использовалась система компьютерной алгебры \textsf{Wolfram Mathematica}. 


Ниже представлен код сценария, использованного для получения значений \(\Psi_{1,2,3}\) при \(\xi\) меняющим значение от 0,1 до 1.

Его также можно найти в репозитории \url{https: //github.com/Fest-Fut/Kursovay}.

\inputminted[linenos,breaklines,fontsize=\scriptsize]{Markdown}{Sun RKF.md}

%\textattachfile{stat-prop_v2.py}{Код в виде прикреплённых данных (только для
%	электронной версии документа).}
